\section{\texorpdfstring{$\sigma$CP}{σCP}-instruments and tensor products}\label{sec:sigmaCP}
Following the standard formalism of von Neumann algebras, \cite{okamura2015measurement} discusses the notion of CP-instruments (\cite{davies1970operational,ozawa1984quantummeasuring}) and how these relate to tensor products of von Neumann algebras.
Here, we provide an analogue of their findings for \cstars{}s.

To fix the setting, recall that a CP-instrument assigns to each measurable subset of a measurable space $(X,\Sigma_X)$ a completely positive endomap of a chosen space, subject to a normalization condition and an adapted form of $\sigma$-additivity. 
In \cite{okamura2015measurement}, the chosen space may be taken to be a
von Neumann algebra $N$, and the endomaps to additionally be \emph{normal}.\footnote{In the reference, this is formalized in a slightly different fashion: the chosen space is the predual, where complete positivity still makes sense. However, since these maps bijectively correspond to normal completely positive maps on the von Neumann algebra, the present formulation is perfectly analogous.}

Here, we adapt this idea to \cstars{}s and \scpu{} maps, and obtain \cref{thm:sigmaCPinstruments}. 
In particular, we prove that $\sigmaCP$-instruments $\Sigma_X \to \sigmaCP(\A,\A)$ are in bijective correspondence with \scpu{} maps $\Linf(X)\kocktensor \A\to \A$. 
We emphasize that the considered maps are not only completely positive, but also $\sigma$-normal, and moreover the domain $\Linf(X)\kocktensor \A$ is itself a \cstars{}.

At first sight, our theorem is a \cstars{} counterpart of \cite[Proposition 3.3]{okamura2015measurement}, which states that every CP-instrument uniquely corresponds to a completely positive map out of a tensor product. 
However, the two are not fully parallel.
As the involved algebras in the reference are von Neumann, one might expect the map to be normal and the tensor product to be a von Neumann algebra, in analogy with \cref{thm:sigmaCPinstruments}.
Instead, the result yields a non-normal map on a
$C^*$-tensor product, so our statement is stronger.

These additional normality properties are not only important from our current perspective. 
In general, the completely positive map of \cite[Proposition 3.3]{okamura2015measurement} need not admit a normal extension to the von Neumann tensor product (\cite[\S V]{okamura2015measurement}), and this obstructs the construction of an associated measuring process (that is, the concrete dynamical realization of an instrument).
To overcome this, the authors introduced the \emph{normal extension property} for instruments, which requires precisely that such an extension exist.
Under this property, \cite[Theorem 3.4]{okamura2015measurement}
establishes a bijective correspondence between CP-instruments and
\emph{normal} completely positive maps from the relevant (von Neumann)
tensor product.

For this reason, \cref{thm:sigmaCPinstruments} captures the spirit of
\cite[Theorem 3.4]{okamura2015measurement}.
Notably, no additional hypothesis in the style of the normal extension
property is needed in our setting: every $\sigmaCP$-instrument
corresponds to a \scpu{} map.
Nonetheless, we make no claim about measuring processes. To the best
of our knowledge, relating our result to them would require a \cstars{}
version of the Stinespring's dilation theorem~\cite[Theorem 4.1]{paulsen2002}, which is not yet
available in the literature.

We expect the findings of this section to be of interest to current
developments in quantum measurement theory and instrument monads, as
they are closely related to the recent preprints
\cite{booth2026quantuminstruments,fritz2026quantuminstrumentmonad}.
In \cite{booth2026quantuminstruments}, the composition of CP-instruments relies on the normal extension property; in contrast,
\cite{fritz2026quantuminstrumentmonad} avoids these extension issues by
restricting to type I von Neumann algebras with separable preduals.
Although a conclusive comparison is beyond the present scope,
\cref{thm:sigmaCPinstruments} (together with
\cref{prop:instrument_composition}) suggests an alternative route in
which these restrictions are not needed, at the price of working with
\cstars{}s and $\sigma$-normal maps.

It is worth noting that, while the aforementioned references treat only measurable spaces, the present discussion applies to arbitrary \bsigma{}s.

%We are now ready to introduce the notion of $\sigmaCP$-instruments.
\begin{definition}
	Let $\Sigma$ be a \bsigma{} and let $\A $ and $\B$ be \cstars{}s. We denote with $\sigmaCP(\A,\B)$ the $\sigma$-normal completely positive maps $\A \to \B$.

	A \newterm{$\sigma$CP-measure} is a map $p\colon \Sigma \to \sigmaCP(\A,\B)$ such that:
	\begin{itemize}
		\item For all sequences $(U_n)_n\subseteq \Sigma$ of pairwise disjoint elements ($U_{n}\wedge U_{m}=\bot$ for all $n\neq m$) and all $a\in \A$, the sequence $(\sum_{n=1}^m p(U_n)(a))_m$ is order convergent in $\B$ and
		      \[\textstyle
			      p\left(\bigvee_n U_n\right)(a) = \sum_n p(U_n)(a)\coloneqq \olim_m^{\B} \left(\sum_{n=1}^m p(U_n)(a)\right)
		      \]
	\end{itemize}
	A \newterm{$\sigmaCP$-instrument} is a {$\sigma$CP-measure} such that $p(\top)(1)=1$.
\end{definition}
Alternatively, the property above may be reformulated by assuming $a\ge 0$ and using $\sup_n$ instead of $\olim_n$.

We now focus on preparatory material for the proof of the main theorem of this section (\cref{thm:sigmaCPinstruments}).

\begin{definition}
	Let $\Sigma$ be a \bsigma{} and let $\B$ be a \cstars{}. A (bounded) \newterm{$\B$-valued measure} is a map $q\colon \Sigma \to \B$ such that
	\begin{itemize}
		\item $q(U)\ge 0$ for all $U\in\Sigma$ and $q(\bot)=0$;
		\item For all sequences $(U_n)_n\subseteq \Sigma$ of pairwise disjoint elements, the sequence $(\sum_{n=1}^m q(U_n))_m$ is order convergent\footnote{Since $q(U_n)$ are positive and $\B$ is $\sigma$-closed, we can alternatively require the existence of $\lambda>0$ such that $\norm{\sum_{n=1}^m q(U_n)}\le \lambda$ for all $m$.} in $\B$
		      \[\textstyle
			      q\left(\bigvee_n U_n\right) = \sum_n q(U_n) \coloneqq \sup_m^{\B} \left(\sum_{n=1}^m q(U_n)\right)
		      \]
	\end{itemize}
\end{definition}
This definition is a generalization of \cite[Definition 2.1]{wright72measures} to the context of \bsigma{}s.
By additionally requiring $q(\top)=1$, one obtains a broader notion of positive operator-valued measures (see \cite[Example 3.2.8 and Definition 4.1.27]{fritz2025categories}).

A crucial remark is that for every \bsigma{} $\Sigma$, there is a valued measure $\Sigma\mapsto \Linfbool(\Sigma)$ denoted by $\chi_{-}$, which sends $U\in \Sigma$ to its abstract ``characteristic function''.
Concretely, $\chi_{U}\colon B(\bbC)\to \Sigma$ is given by $\chi_{U}(\lbrace 1\rbrace)=U$ and $\chi_{U}(\lbrace 0\rbrace)=\neg U$ (recall the description given in \cref{sec:preliminaries}).

In particular, for any sequence $(U_n)$,
\[
	\chi_{\bigvee_n U_n}=\sup_m\chi_{\bigvee_{n=1}^m U_n}\qquad\text{and}\qquad\chi_{\bigwedge_n U_n}= \inf_m \chi_{\bigwedge_{n=1}^m U_n}
\]
while for the bottom and top elements we have $\chi_{\bot}=0$ and $\chi_{\top}=1$.

For the categorically inclined, $\chi_{-}$ is the unit of the adjoint equivalence between $\cSigmaC$ and the category of \bsigma{}s with \bsigmahom{}s.
\begin{theorem}\label{thm:VM_sigmacp}
	Let $\Sigma$ be a \bsigma{} and $\B$ a \cstars{}.
	Then every $\B$-valued measure $q$ uniquely extends to a $\sigma$-normal completely positive map $\tilde{q}\colon \Linfbool(\Sigma)\to \B$ in the sense of the following commutative diagram:
	\[
		\xymatrix{\Sigma\ar[rr]^{\chi_{-}}\ar[rd]_{q}&&\Linfbool(\Sigma)\ar[ld]^{\widetilde{q}}\\
		&\B&
		}
	\]
\end{theorem}
\begin{proof}[Sketch of the proof]
	The argument is a simple rewriting of the proof of \cite[Theorem 4.1.30]{fritz2025categories}.
	Let us offer some details for the sake of clarity.
	If $\Sigma$ is an honest $\sigma$-algebra, i.e.\ there exists a measurable space $(X,\Sigma_X)$ with $\Sigma_X\cong \Sigma$, then this result was shown in \cite[Section 2]{wright72measures} and basically follows by rewriting the construction of the Lebesgue integral.
	For a general \bsigma{} $\Sigma$, by Loomis--Sikorski representation theorem \cite[Theorem 29.1]{sikorski69boolean}, there exists a measurable space $(X,\Sigma_X)$ such that $\Sigma\cong \Sigma_X/I$, with $I$ a $\sigma$-ideal.
	In particular, $\Linfbool(\Sigma)\cong \Linf(X)/\mathcal{I}$, where $\mathcal{I}$ is the set of bounded measurable functions whose support lies in $I$ (\cite[Proposition 4.1.22]{fritz2025categories}).
	We now consider
	\[
		\xymatrix{\Sigma_{X}
		\ar[rrrr]^{\chi_-}
		\ar@{->>}[dr]_{r}
		&&&&
		\Linf(X)
		\ar@{->>}[dl]^{\widetilde{\chi_{-} r}}
		\ar@/^2.5pc/[ddll]^{\widetilde{qr}}
		\\
		&
		\Sigma
		\ar[rr]^{\chi_-}
		\ar[dr]_{q}
		&&
		\Linf(X)/\mathcal{I}
		\ar@{-->}[dl]^{\widetilde{q}}
		\\
		&&
		\A}
	\]
	where $\widetilde{qr}$ exists from the previous step.
	We claim that the dashed map exists by the quotient universal property.
	Indeed, for any self-adjoint $g\in \mathcal{I}$, we have
	\[
		-\norm{g}\chi_{\operatorname{supp}(g)} \le g\le \norm{g} \chi_{\operatorname{supp}(g)}
	\] with $\operatorname{supp}(g)\in I$ the support of $g$ (multiply the standard inequality $-\norm{g}\le g\le \norm{g}$ by $\chi_{\operatorname{supp}(g)}$).
	Together with $\widetilde{qr}(\chi_{\operatorname{supp}(g)})=qr({\operatorname{supp}(g)})= 0$, this implies
	\[ 0 =-\norm{g}\cdot 0 = \widetilde{qr}(-\norm{g}\chi_{\operatorname{supp}(g)}) \le \widetilde{qr}(g)\le \widetilde{qr}(\norm{g} \chi_{\operatorname{supp}(g)}) = \norm{g} \cdot 0 =0,
	\]
	and therefore $\widetilde{qr}(g)=0$.
\end{proof}
For the next preliminary result, we recall the following.
\begin{definition}
	In a \cstar{} $\A$, a \newterm{projection} $p$ is an element of $\A$ such that $p^*=p=p^2$.

	A \newterm{simple element} $s$ is a linear combination of projections: $s=\sum_{k=1}^m \lambda_k p_k$.
\end{definition}
In $\Linf(X)$, a simple element is given by a simple function because projections are characteristic functions of measurable subsets.
Similarly, in $\Linfbool(\Sigma)$, projections are in correspondence with the elements of $\Sigma$ via $\chi_{-}\colon \Sigma \to \Linfbool(\Sigma)$ (\cite[Lemma 4.1.26]{fritz2025categories}).

\begin{lemma}\label{lem:technical_simple_els}
	Let $\Sigma$ be a \bsigma{} and let $\A$ be a \cstars{}.
	\begin{enumerate}
		\item\label{it:simple_els_dense} The linear span of elements $s\otimes a$, where $a\in \A$ and $s$ is a simple element in $\Linfbool(\Sigma)$, is dense in $\Linfbool(\Sigma)\otimes \A$.
		      Moreover, the self-adjoint elements of this span is dense in the set of self-adjoint elements of $\Linfbool(\Sigma)\otimes \A$.
		\item\label{it:simple_els_positivity} $\sum_{k=1}^m s_k \otimes a_k$ is positive if and only there exist some $n$ and some $t_{\ell}\ge 0$ and $b_{\ell}\ge 0$ for every $\ell \in \lbrace 1, \dots , n\rbrace$, such that
		      \[
			      \sum_{k=1}^m s_k \otimes a_k = \sum_{\ell=1}^n t_{\ell} \otimes b_{\ell}.
		      \]
		      Moreover, $t_{\ell}$ can be taken to be disjoint projections: $t_{\ell}\cdot t_{\ell'}=0$ for all distinct $\ell,\ell'$.
	\end{enumerate}
\end{lemma}
\begin{proof}
	Item~\ref{it:simple_els_dense} follows from the fact that linear combinations of $f\otimes a$ are dense by definition of the tensor product (since it is the norm-closure of the algebraic tensor product), and the fact that simple elements are dense in $\Linfbool(\Sigma)$ by \cite[Corollary 3.4.4]{fritz2025categories}.
	The restriction to self-adjoint elements is analogous.

	Regarding Item~\ref{it:simple_els_positivity}, let us assume that $\sum_{k=1}^m s_k \otimes a_k\ge 0$.
	Let $s_k= \sum_{j} \lambda_{j,k} p_{j,k}$ with $\lambda_{j,k}\in\bbC$ and $p_{j,k}$ projections.
	By standard properties of the algebraic tensor product,
	\[\textstyle
		\sum_k s_k \otimes a_k =  \sum_k \Big(\sum_j  \lambda_{j,k} p_{j,k}\Big)\otimes a_k = \sum_{j,k} p_{j,k}\otimes (\lambda_{j,k}a_k)
	\]
	Hence, we can assume without loss of generality that our element is written as $\sum_{k=1}^m p_k \otimes a_k$ with $p_k$ projections.
	For the next step, we introduce some auxiliary notation. Let $p_k^{\top}\coloneqq p_k$ and $p_k^{\bot}\coloneqq 1-p_k$.
	Now, for any Boolean sequence $I=(i_1, \dots, i_m) \in \lbrace \bot,\top\rbrace^m$, we define
	\[
		p^I \coloneqq (p_1)^{i_1} \cdot (p_2)^{i_2} \cdots (p_m)^{i_m}, \qquad \text{and}\qquad a^I \coloneqq \sum_{j \,\mid\, i_j=\top} a_j
	\]
	By a simple induction reasoning,
	\[
		\sum_{k=1}^m p_k \otimes a_k = \sum_{I \in \lbrace \bot,\top\rbrace^m} p^I \otimes a^I,
	\]
	and moreover $p^I \cdot p^J=0$ for any distinct $I,J$ because $p_k^{\top}p_k^{\bot} =0$ for any $k$.
	By restricting to those $I\in \lbrace \bot,\top\rbrace^m$ such that $p^I\neq 0$, we can set $t_{\ell}=p^I$ and $a^I=b^{\ell}$ and obtain the formula described in the statement.
	However, we have yet to show that $a^I$ is positive (while positivity of $p^I$ is immediate since it is a projection).
	To this end, we recall the well-known fact that, for any compact Hausdorff space $Y$, $C(Y)\otimes \A \cong C(Y,\A)$, where $C(Y,\A)$ is the \cstar{} given by continuous maps $Y\to \A$ where operations are defined pointwise (see, e.g., \cite[Corollary T.6.17]{weggeolsen1993Ktheory}).
	From this description, one notes that $f\in C(Y,\A)$ is positive if and only if $f(y)\ge 0$ for any $y\in Y$.
	By taking $Y$ the Gelfand spectrum of $\Linfbool(\Sigma)$ and $U_I $ the subset of $Y$ satisfying $p^I = \chi_{U_I}$,\footnote{$p^I$ interpreted as a function $Y\to\bbC$ has image contained in $\lbrace 0,1\rbrace$ because it is a projection.} for any element $u\in U^I$,
	\[
		p^J(u)=p^J(u)p^I(u)= 0
	\] for any $J\neq I$, because $p^J p^I=0$.
	Moreover, $U^I$ is nonempty because we restricted to those $I$ such that $p^I\neq 0$.
	It follows that $a^I$ is exactly the image of $u$ via the continuous map $\sum_{I \in \lbrace \bot,\top\rbrace^m} a^I \cdot p^I(-) \colon Y \to \A$, which is positive by assumption.
	Finally, $a^I$ is positive, and the statement is shown.
\end{proof}

\begin{theorem}\label{thm:sigmaCPinstruments}
	Let $\Sigma$ be a \bsigma{} and let $\A$ and $\B$ be \cstars{}s.
	There is a bijective correspondence between
	\begin{enumerate}
		\item\label{it:sigmaCPinstrument} $\sigmaCP$-instruments $\Sigma \to \sigmaCP(\A,\B)$;
		\item\label{it:scpu_tensor} \scpu{} maps $\Linfbool(\Sigma)\kocktensor \A \to \B$.
	\end{enumerate}
\end{theorem}
\begin{proof}
	We define a map $p_{-}$ from Item~\ref{it:scpu_tensor} to Item~\ref{it:sigmaCPinstrument} by sending a \scpu{} map $\phi\colon \Linfbool(\Sigma)\kocktensor \A \to \B$ to the $\sigmaCP$-instrument $p_{\phi}\colon U \mapsto (a\mapsto \phi(\chi_U \otimes a))$.
	This is indeed a $\sigmaCP$-measure: for every sequence $(U_n)\subseteq \Sigma$ of pairwise disjoint elements,
	\begin{align*}
		\textstyle p_{\phi}\left(\bigvee_n U_n\right)(a) & \textstyle =\phi(\chi_{\bigvee_n U_n} \otimes a)                                                                                                  \\
		                                                 & \textstyle = \phi \left(\left(\sup_m \sum_{n=1}^m\chi_{U_n} \right)\otimes a\right)   \tag{$\chi_{\bigvee_n U_n}=\sup_m \sum_{n=1}^m \chi_{U_n}$} \\
		                                                 & \textstyle = \olim_m \phi\left(\left(\sum_{n=1}^m\chi_{U_n} \right) \otimes a \right) \tag{\cref{prop:universal_bisigma}}                         \\
		                                                 & \textstyle = \olim_m \sum_{n=1}^m \phi(\chi_{U_n}\otimes a)                           \tag{linearity}                                             \\
		                                                 & \textstyle = \sum_{n} p_{\phi} (U_n)(a)
	\end{align*}
	Normalization is obvious: $p_{\phi}(\top)(1)=\phi(\chi_{\top}\otimes 1)= \phi(1\otimes 1)=1$.

	Injectivity of $p_{-}$ is also straightforward.
	If $p_{\phi}=p_{\psi}$, then $\phi(\chi_U\otimes a)=\psi(\chi_U\otimes a)$ for all $a\in \A$ and all $U\in\Sigma$. Since for every $f\in \Linfbool(\Sigma)$ one has $f=\olim_n s_n$, where $s_n$ are simple elements (\cite[Corollary 3.4.4]{fritz2025categories}), this is sufficient to conclude that $\phi$ and $\psi$ are equal on all simple tensors by bi-$\sigma$-normality (\cref{prop:universal_bisigma}), and therefore they must coincide.

	We are left to show surjectivity of $p_{-}$.
	We first note that, for any positive $a\in \A$, a $\sigmaCP$-instrument $p$ yields a $\B$-valued measure $q_a(U)\coloneqq p(U)(a)$.
	The $\sigma$-additivity of $q_a$ is immediate from that of $p$, so we only need to check $q_a(U)\ge 0$ and $q_a(\bot)=0$.
	The first one follows from the fact that $a$ is positive and $p(U)$ is completely positive.
	The second one is a consequence of $q_a(U\vee \bot)=q_a(U)+q_a(\bot)$ for all $U\in \Sigma$ by $\sigma$-additivity of $q_a$: indeed, this entails $q_a(\bot)=2q_a(\bot)$, and so $q_a(\bot)=0$.
	By \cref{thm:VM_sigmacp}, we obtain a $\sigma$-normal cp map $\Linfbool(\Sigma)\to \B$ extending $q_a$.
	Moreover, whenever $f=\olim_n s_n$ with $s_n=\sum_\ell \lambda_{\ell,n} \chi_{U_{\ell,n}}$ simple elements, one has $\widetilde{q_a} (f)= \olim_n \sum_\ell \lambda_{\ell,n} p({U_{\ell,n}})(a)$ by $\sigma$-normality of $\widetilde{q_a}$.
	Overall, we have obtained a map
	\[
		\phi \colon   \xymatrix@R=.1pt{
		\Linfbool(\Sigma)\times \A_+ \ar[r]& \B \\
		(f,a)\ar@{|->}[r]
		& \widetilde{q_a}(f)
		}
	\]
	where $\A_+$ denotes the positive elements of $\A$.
	By linearity, we can extend the whole discussion to all elements of $\A$.
	Explicitly, given $a=\sum_i \mu_i a_i = \sum_j \nu_j a_j'$ finite sums with $a_i,a_j'\ge 0$ and $\mu_i,\nu_j \in \bbC$, we have
	\begin{align*}
		\textstyle\sum_i\mu_i \phi(f,a_i) & =\textstyle \sum_i \mu_i \left(\olim_n \sum_{\ell} \lambda_{\ell,n} p(U_{\ell,n})(a_i)\right)          \\
		                                 & =\textstyle \olim_n \sum_{\ell} \lambda_{\ell,n}p(U_{\ell,n}) \left(\sum_i \mu_i a_i\right) \tag{\dag} \\
		                                 & =\textstyle\olim_n \sum_{\ell} \lambda_{\ell,n}p(U_{\ell,n}) \left(\sum_j \nu_j a_j'\right)            \\
		                                 & =\textstyle \sum_j \nu_j \olim_n  \sum_{\ell} \lambda_{\ell,n} p(U_{\ell,n})(a_j') \tag{\ddag}          \\
		                                 & =\textstyle\sum_j \nu_j \phi(f,a_j')
	\end{align*}
	where the two equalities $(\dag),(\ddag)$ hold by linearity of order limits and $p(U)$, for all $U \in \Sigma$.
	% In particular,
	% \[
	% \phi(f,a) = \olim_n \sum_{\ell} \lambda_{\ell,n} p(U_{\ell,n})(a)
	% \]
	% holds for all elements $a\in \A$ and all choices of $f= \olim_n \sum_\ell \lambda_{\ell,n} \chi_{U_{\ell,n}}$.
	Finally, we obtain a linear map $\Linfbool(\Sigma)\odot \A \to \B $, which we again call $\phi$ for the sake of simplicity, where $\odot$ denotes the algebraic tensor product.

	By \cref{lem:technical_simple_els}\ref{it:simple_els_positivity}, we see that this map restricted to the elements $\sum_{k=1}^m s_k\otimes a_k$ is positive, and therefore, by \cite[Proposition 2.1]{paulsen2002} and its proof, it is bounded and satisfies $\norm{\phi(b)}\le \norm{b}$ for all self-adjoint $b$.
	Finally, $\phi$ can be uniquely extended to the norm-completion of its domain, which coincides with $\Linfbool(\Sigma)\otimes \A$ by \cref{lem:technical_simple_els}\ref{it:simple_els_dense}.
	Moreover, such an extension must be self-adjoint, and contractive on self-adjoint elements, by the second sentence of \cref{lem:technical_simple_els}\ref{it:simple_els_dense}.
	We can now apply \cite[Corollary 3.25]{fritz2023involutive} to conclude that this map is positive.

	By reapplying the same strategy to the $\sigmaCP$-instrument $\Sigma \to \sigmaCP(\A\otimes M_n , \B\otimes M_n)$ sending $U$ to $p(U)\otimes \id_{M_n}$, we obtain positivity of $\Linfbool(\Sigma)\otimes \A\otimes M_n \to \B \otimes M_n$.
	Arbitrariness of $n$ ensures complete positivity.

	We now check bi-$\sigma$-normality of $\phi\colon \Linfbool(\Sigma)\otimes \A\to \B$.
	On $\Linfbool(\Sigma)$, $\sigma$-normality follows from construction: given $a\in \A$ and $a=\sum_i \mu_i a_i$ with $a_i\ge 0$ and $\mu_i \in \bbC$,
	\begin{align*}
		\phi(\olim_k f_k \otimes a) & =\sum_i \mu_i \phi(\olim_k f_k \otimes a_i) = \sum_i \mu_i q_{a_i}(\olim_k f_k)  \\
		                            & = \sum_i \mu_i \olim_k q_{a_i}(f_k) = \olim_k \sum_i \mu_i \phi(f_k \otimes a_i) \\
		                            & = \olim_k \phi(f_k\otimes a)
	\end{align*}
	(see also \cite[Theorem 4.1.7]{fritz2025categories}).

	The other side is more delicate.
	Let us restrict to the case of $f\in \Linfbool(\Sigma)$ positive, and let $(a_n)\subseteq \A $ be a monotone decreasing sequence such that $\inf_n a_n =0$; in particular, $a_n\ge 0$ for all $n$.
	Then, $\phi(f\otimes a_n)\le \phi(f\otimes a_{n-1})$, and we wish to show that $\inf_n \phi(f\otimes a_n) =0$.
	To this end, note that $0 \le \phi(f\otimes a_n)\le\phi(\norm{f}\otimes a_n)$ by positivity of $\phi$ and the standard inequality $0\le f\le \norm{f}$.
	Consequently,
	\[\textstyle
		0 \le \inf_n \phi(f\otimes a_n) \le \inf_n \phi(\norm{f}\otimes a_n) \overset{(\dag)}{=} \inf_n \norm{f} p(\top)(a_n) = \norm{f}p(\top)(\inf_n a_n) = 0
	\]
	where the equality $(\dag)$ follows from definition of $\phi$, while the next equality holds by $\sigma$-normality of $p(\top)$. Therefore, $\inf_n \phi(f\otimes a_n)=0=\phi(f\otimes \inf_n a_n)$, as wanted. By linearity, $\phi(f\otimes -)$ also preserves any suprema of norm-bounded monotone increasing sequences (cf.\ \cite[Remark 3.2.5]{fritz2025categories}), and so it preserves also any order limit (as mentioned in \cref{sec:preliminaries}).
	Since every element of a \cstar{} is a linear combination of positive elements, $\phi(f \otimes \olim_m a_m) = \olim_m \phi(f\otimes a_m)$ holds for all $f\in \Linfbool(\Sigma)$.

	Finally, bi-$\sigma$-normality of $\phi$ is ensured, and surjectivity of $p_{-}$ now follows from \cref{prop:universal_bisigma}.
\end{proof}

We close with the following composition of $\sigmaCP$-instruments,
a first step toward a variant for \cstars{}s of the program of quantum
instrument monads~\cite{booth2026quantuminstruments,fritz2026quantuminstrumentmonad}.

\begin{proposition}\label{prop:instrument_composition}
		Let $\Sigma$ and $\Sigma'$ be \bsigma{}s and let $\A$, $\B$, $\C$ be \cstars{}s.
		$\sigmaCP$-instruments $p\colon \Sigma \to \sigmaCP(\A,\B)$ and $p'\colon \Sigma' \to \sigmaCP(\B,\C)$ compose to a $\sigmaCP$-instrument
		\[
			p' \circ p \colon \Sigma \otimes \Sigma' \longrightarrow \sigmaCP(\A,\C)
		\]
		whose outcomes are indexed by the universal tensor product of \bsigma{}s.
\end{proposition}
\begin{proof}
		By \cref{thm:sigmaCPinstruments}, $p$ and $p'$ correspond to \scpu{} maps $\phi\colon \Linfbool(\Sigma)\kocktensor \A \to \B$ and $\phi'\colon \Linfbool(\Sigma')\kocktensor \B \to \C$.
		Since $\SigmaCcpu$ is symmetric monoidal (\cref{thm:tensor_scpu}), the composite
		\[
			\Linfbool(\Sigma')\kocktensor \Linfbool(\Sigma)\kocktensor \A
			\xrightarrow{\ \id \kocktensor \phi\ }
			\Linfbool(\Sigma')\kocktensor \B
			\xrightarrow{\ \phi'\ }
			\C
		\]
		is again a \scpu{} map.
		By \cref{prop:tensor_commutative} and \cref{rem:tensor_bool}, $\Linfbool(\Sigma')\kocktensor \Linfbool(\Sigma)\cong \Linfbool(\Sigma\otimes \Sigma')$, where $\otimes$ denotes the universal tensor product of \bsigma{}s~\cite[Proposition 2.5.1]{fritz2025categories}.
		Applying \cref{thm:sigmaCPinstruments} in the other direction yields the desired instrument.
\end{proof}




